\section*{Chapitre \ref{ch:nw:the-asia-pacific-map} --- La carte Asie-Pacifique}

\begin{solution}{exo:nw:the-asia-pacific-map:1}
Sur \qty{6296.4}{\kilo\metre}~: \qty{21.00}{\milli\second} dans le vide et \qty{30.71}{\milli\second} dans une fibre droite, dans un sens.
\end{solution}

\begin{solution}{exo:nw:the-asia-pacific-map:2}
Neuf copies plus tard~: $9 \times 10 = \qty{90}{\micro\second}$ (105 contre \qty{15}{\micro\second} sur le premier port du premier serveur, sans gigue).
\end{solution}

\begin{solution}{exo:nw:the-asia-pacific-map:3}
Entre Tokyo et Hong Kong le plancher est \qty{9.6}{\milli\second}~; le situer à l'intérieur de sa ville le décale de quelques dizaines de microsecondes, soit moins de 1~\%.
À l'intérieur d'un centre de données, la concurrence se joue à l'échelle de la nanoseconde ou de la microseconde, si bien que la position exacte d'un bâtiment et, à l'intérieur de celui-ci, la longueur des câbles
décident de la course.
\end{solution}

\begin{solution}{exo:nw:the-asia-pacific-map:4}
L'expéditeur envoie toujours une copie après l'autre, au même coût, si bien que la première et la dernière copie restent aussi éloignées qu'avant~; ce qui change, c'est le membre qui reçoit chaque copie.
Un ordre fixe donne au même membre la première copie à chaque tick~; un ordre aléatoire fait que chaque membre reçoit la première copie à certains ticks,
et la dernière à d'autres, si bien que personne n'a d'avantage persistant (corrélation de rang proche de zéro).
\end{solution}

\begin{solution}{exo:nw:the-asia-pacific-map:5}
Ses ports n'ont que deux membres au lieu de dix, si bien que, même en dernière position dans l'ordre du centre, ses membres reçoivent le tick après
$5 \times 5 + 10$ à $5 \times 5 + 20$, soit de 35 à \qty{45}{\micro\second}, avant la plupart des membres des serveurs chargés, quel que soit le mélange. Un
mécanisme de randomisation supprime l'avantage d'être connecté le premier, pas l'avantage d'être sur un serveur moins chargé~: c'est pour cette raison que
l'ordonnance de la SEBI traite les serveurs secondaires comme un problème distinct.
\end{solution}

\begin{solution}{exo:nw:the-asia-pacific-map:6}
Multicast pour la stratégie~: chaque récepteur installé en \omterm{def:nw:colocation-products-and-how-they-are-sold:colo}{colocation} reçoit la même copie au même instant, à la réplication dans les commutateurs et à son propre câble près.
TCP pour la reprise et pour les lieux que le multicast n'atteint pas (un site de reprise après sinistre, un bureau distant)~; les questions fréquentes
indiquent que le moteur du centre de données secondaire est joint en TCP.
\end{solution}

\begin{solution}{exo:nw:the-asia-pacific-map:7}
\texttt{first\_vs\_median(m)} donne 15.5, 28.5, 50.4 et \qty{100.5}{\micro\second} pour $m = 2, 5, 10, 20$, contre $c\,m/2 = 10$, 25, 50 et
\qty{100}{\micro\second}. La proposition capture la croissance~; le léger excès pour les petites valeurs de $m$ vient de ce que le serveur du membre médian se trouve plus tard
dans l'ordre du centre, soit quelques copies de \qty{5}{\micro\second}.
\end{solution}

\begin{solution}{exo:nw:the-asia-pacific-map:8}
Des câbles de même longueur égalisent le chemin entre le serveur et chaque membre, pas l'instant auquel le serveur envoie sa copie à chaque membre. Avec une
diffusion TCP séquentielle, le dixième membre d'un port recevait chaque tick neuf envois après le premier, quelle que soit la longueur de son câble.
\end{solution}

\begin{solution}{pb:nw:the-asia-pacific-map:1}
\textbf{Partie I.}
\begin{enumerate}
\item Premier après $5 + 10 = \qty{15}{\micro\second}$, dernier après $5 + 100 = \qty{105}{\micro\second}$ (plus la gigue).
\item Les cinquième et sixième membres reçoivent après 55 et \qty{65}{\micro\second}~: \qty{60}{\micro\second} en moyenne.
\item Tout est décalé de $3 \times 5 = \qty{15}{\micro\second}$~: 30 pour le premier membre, 120 pour le dernier.
\item Premier~: le premier membre de chaque port du premier serveur, à \qty{15}{\micro\second}~; dernier~: les dixièmes membres du quatrième serveur, à
  \qty{120}{\micro\second}.
\end{enumerate}
\textbf{Partie II.}
\begin{enumerate}[resume]
\item 15.5, 28.5, 50.4 et \qty{100.5}{\micro\second}.
\item Proche de $c\,m/2$~: 10, 25, 50, 100~; l'excès vient de l'ordre du centre.
\item Dans tous les cas~: 100~\% dans le modèle, car l'ordre ne change jamais et \qty{50}{\micro\second} écrase le bruit de
  \qty{5}{\micro\second}.
\item Décisif~: les courses sur cotations périmées se jouent à la microseconde près ou moins, donc un membre \qty{50}{\micro\second} en avance gagne presque toutes celles
  auxquelles il participe.
\end{enumerate}
\textbf{Partie III.}
\begin{enumerate}[resume]
\item L'écart reste d'environ \qty{110}{\micro\second}~; la corrélation de rang passe de 1.00 à 0.07~; le taux de victoire persistant ne tombe qu'à 88~\%,
  car le serveur secondaire, moins chargé, garde ses membres en avance.
\item Le taux de victoire tombe à 51~\%~: personne n'est systématiquement premier.
\item L'écart tombe à la gigue (\qty{8.4}{\micro\second} aux extrêmes, \qty{0.7}{\micro\second} du premier au médian) et le taux de
  victoire à 50~\%.
\item Multicast~: il supprime à la fois la persistance et l'écart~; les deux autres ne font que redistribuer un écart qui subsiste.
\end{enumerate}
\textbf{Partie IV.}
\begin{enumerate}[resume]
\item \emph{Résultat nommé}~: avec dix membres par port, le premier membre reçoit chaque tick environ \qty{50}{\micro\second} avant le membre médian
  en unicast séquentiel, et environ \qty{0.7}{\micro\second} avant en multicast~; en unicast avec un ordre fixe, il remporte chaque course
  contre le membre médian.
\item Faits~: l'ordre séquentiel à deux niveaux, trois ports par serveur, l'ordre de connexion comme ordre d'envoi, un serveur secondaire peu chargé.
  Hypothèses~: les coûts par copie, la gigue, dix membres par port, le modèle de temps de réaction.
\item Horodater les mêmes messages sur plusieurs connexions, ports et serveurs, avec des horloges synchronisées (chapitres~4 et~5)~; calculer l'écart des
  instants de réception par message et la corrélation de rang entre messages.
\item Parce que ce qui compte est l'instant auquel la donnée atteint chaque membre, lequel dépend de toute la chaîne (ordre d'envoi, charge
  des serveurs, câbles), et non sur un seul de ses maillons.
\item Comme exigence, la latence de communication entre les zones de \omterm{def:nw:colocation-products-and-how-they-are-sold:colo}{colocation} dédiées et les processeurs frontaux doit être cohérente.
\item Flux WebSocket issus de plateformes crypto (chapitre~20), relais hébergés sur un cloud public (chapitre~16), et toute diffusion en éventail d'un fournisseur sur TCP.
\item Un chemin plus court relève de la géographie ou de l'ingénierie, ouvert à quiconque le paie~; un ordre d'envoi est un choix de l'émetteur,
  que la plateforme équitable doit rendre indépendant de l'identité du destinataire.
\item Parce que le flux envoyait chaque tick aux membres l'un après l'autre, dans l'ordre de leur connexion, si bien que le premier connecté voyait chaque tick
  en premier.
\end{enumerate}
\end{solution}

\begin{solution}{iq:nw:the-asia-pacific-map:1}
Unicast envoie une copie par destinataire sur sa propre connexion, l'une après l'autre~; multicast envoie une seule copie à une adresse de groupe que le réseau
réplique. En \omterm{def:nw:colocation-products-and-how-they-are-sold:colo}{colocation}, le multicast donne à chaque récepteur la même copie au même instant et ne coûte à la plateforme qu'un envoi par message,
quel que soit le nombre de récepteurs.
\iqlookfor{Envois séquentiels et leur ordre~; réplication dans les commutateurs~; perte et reprise sur un canal unicast séparé.}
\end{solution}

\begin{solution}{iq:nw:the-asia-pacific-map:2}
Tokyo--Hong Kong environ \qty{2900}{\kilo\metre} (\qty{9.6}{\milli\second} dans le vide, \qty{14}{\milli\second} en fibre), Hong Kong--Singapour environ
\qty{2600}{\kilo\metre} (\qty{8.6}{\milli\second}), Tokyo--Singapour environ \qty{5300}{\kilo\metre} (\qty{17.7}{\milli\second}). Les stratégies inter-marchés
y fonctionnent à des horizons de l'ordre de la milliseconde et doivent choisir un marché de rattachement pour chaque moteur.
\iqlookfor{Bon ordre de grandeur~; aller simple contre aller-retour~; choix de l'emplacement de chaque moteur.}
\end{solution}

\begin{solution}{iq:nw:the-asia-pacific-map:3}
Il faut d'abord vérifier la capture (les deux horodatées par un matériel synchronisé au même point)~; puis comparer les serveurs, ports et heures de connexion des deux connexions, et voir si le retard est constant (un ordre d'envoi),
s'il croît par rafales (un récepteur lent ou une file d'attente) ou apparaît après une reconnexion~; il faut ensuite se reconnecter dans un autre ordre
et vérifier si le retard se déplace.
\iqlookfor{La mesure avant la théorie~; file d'attente contre ordre d'envoi~; une expérience qui les sépare.}
\end{solution}

\begin{solution}{iq:nw:the-asia-pacific-map:4}
Que chaque destinataire d'un même service reçoive chaque message au même instant, à un écart près qui ne dépend pas de son identité. On le teste en horodatant
les mêmes messages sur plusieurs connexions et en vérifiant l'écart et si l'ordre des destinataires persiste.
\iqlookfor{Les deux volets~: écart et persistance~; les rapports de latence publiés par la plateforme~; des conditions égales au sein d'un même niveau de service.}
\end{solution}

\begin{solution}{iq:nw:the-asia-pacific-map:5}
Celui vers qui le serveur écrit en premier~: en général l'ordre de connexion ou d'une liste interne, éventuellement différent à chaque mise à jour~; le dernier client
attend toutes les écritures qui le précèdent, et un client lent peut retarder les autres si le serveur écrit de façon synchrone.
\iqlookfor{Le coût du fan-out croît avec le nombre de clients~; ordre d'envoi~; files par client et consommateurs lents.}
\end{solution}

\begin{solution}{iq:nw:the-asia-pacific-map:6}
Multicast en \omterm{def:nw:colocation-products-and-how-they-are-sold:colo}{colocation}, avec des câbles de même longueur depuis un point de réplication unique~; un niveau de service par produit, assorti de conditions publiées~; un
canal de reprise unicast qui ne soit pas plus rapide que le multicast~; des horodatages par message à chaque transmission entre membres, et des statistiques publiées de
l'écart et de la persistance de l'ordre.
\iqlookfor{Conception et preuves ensemble~; le chemin de reprise comme avantage caché~; mesure indépendante.}
\end{solution}
